\documentclass{article}
\usepackage{parskip}
\usepackage{booktabs}
\usepackage{geometry}
\geometry{margin=1in}
\title{Practice Problem 2.6}
\author{CSAPP}
\date{}
\begin{document}
\maketitle

\begin{flushleft}
    Using \texttt{show\_int} and \texttt{show\_float}, we determine that the
    integer 2607352 has hexadecimal representation \texttt{0x0027C8F8}, while
    the floating-point number 3510593.0 has hexadecimal representation
    \texttt{0x4A1F23E0}.
\end{flushleft}

\bigskip

\hspace{1em}\section*{\small Solution A}
\begin{center}
    \footnotesize
    \begin{tabular}{l l l}
        \toprule
        \textbf{Hex}    & \texttt{0x0027C8F8} & \texttt{0x4A1F23E0} \\
        \midrule
        \textbf{Binary} & \texttt{0000~0000~0010~0111~1100~1000~1111~1000} & \texttt{0100~1010~0001~1111~0010~0011~1110~0000} \\
        \bottomrule
    \end{tabular}
\end{center}

\bigskip

\section*{\small Solution B}
\begin{center}
    \footnotesize
    \begin{tabular}{l l}
        \toprule
        \textbf{Binary 1} & \texttt{0000~0010~0111~1100~1000~1111~1000} \\
        \textbf{Binary 2} & \texttt{0100~1010~0001~1111~0010~0011~1110} \\
        \textbf{Match}    & \texttt{!x!!~x!!!~!xx!~!!xx~x!x!~xx!!~!xx!} \\
        \bottomrule
    \end{tabular}
\end{center}

\bigskip

\section*{\small Solution C}
\begin{flushleft}
    \footnotesize
    The non-matching parts are the bits that differ after alignment.
    Specifically:
    \begin{itemize}
        \item The leading \texttt{0100~1010} of \texttt{0x4A1F23E0} 
              (sign bit and exponent field of the float) has no corresponding 
              match in the integer representation.
        \item The trailing \texttt{0000} of \texttt{0x4A1F23E0} 
              (least significant bits of the float) does not appear in the 
              integer.
        \item The leading \texttt{0000~0000} of \texttt{0x0027C8F8} 
              (most significant zero bits of the integer) fall outside the 
              float's mantissa range.
    \end{itemize}
\end{flushleft}

\end{document}